\documentclass[a4paper,11pt]{article}
	\usepackage[utf8]{inputenc}
	\usepackage[italian]{babel}
	\usepackage{hyperref}	%Consente l'inserimento di \url
	\usepackage{booktabs}	%Utilità di abbellimento tabelle
	\usepackage{longtable}
	\usepackage{tabularx}
	%\usepackage{widetable}
	\usepackage{array}
	\usepackage{listings}
	\usepackage{graphicx}
	\usepackage{caption}
	\usepackage{fancyhdr}
	\newenvironment{fixpic}{}{} % [1]
	\usepackage[a4paper,top=3cm,bottom=3cm,left=2.5cm,right=2.5cm]{geometry}
	%******
	\usepackage{makeidx}
	\usepackage{textcomp}
	\usepackage{multirow}
	\usepackage{rotfloat}
	\usepackage{lastpage}
	\usepackage{array}
	\usepackage{float}
	% *************************************
	% QUI CODICE PER \SUBSUBSUBSECTION
	\usepackage{titlesec}
	\titleclass{\subsubsubsection}{straight}[\subsection]
	
	\newcounter{subsubsubsection}[subsubsection]
	\renewcommand\thesubsubsubsection{\thesubsubsection.\arabic{subsubsubsection}}
	\renewcommand\theparagraph{\thesubsubsubsection.\arabic{paragraph}} % optional; useful if paragraphs are to be numbered
	
	\titleformat{\subsubsubsection}
	  {\normalfont\normalsize\bfseries}{\thesubsubsubsection}{1em}{}
	\titlespacing*{\subsubsubsection}
	{0pt}{3.25ex plus 1ex minus .2ex}{1.5ex plus .2ex}
	
	\makeatletter
	\renewcommand\paragraph{\@startsection{paragraph}{5}{\z@}%
	  {3.25ex \@plus1ex \@minus.2ex}%
	  {-1em}%
	  {\normalfont\normalsize\bfseries}}
	\renewcommand\subparagraph{\@startsection{subparagraph}{6}{\parindent}%
	  {3.25ex \@plus1ex \@minus .2ex}%
	  {-1em}%
	  {\normalfont\normalsize\bfseries}}
	\def\toclevel@subsubsubsection{4}
	\def\toclevel@paragraph{5}
	\def\toclevel@paragraph{6}
	\def\l@subsubsubsection{\@dottedtocline{4}{7em}{4em}}
	\def\l@paragraph{\@dottedtocline{5}{10em}{5em}}
	\def\l@subparagraph{\@dottedtocline{6}{14em}{6em}}
	\makeatother
	
	\setcounter{secnumdepth}{4}
	\setcounter{tocdepth}{4}
	%FINE \SUBSUBSUBSECTION
	%****************************************
	%STYLE PER INSERIMENTO DEL CODICE
	\lstdefinestyle{style1}{
	  belowcaptionskip=1\baselineskip,
	  breaklines=true,
	  frame=L,
	  xleftmargin=\parindent,
	  language=Pascal,
	  showstringspaces=false,
	  basicstyle=\footnotesize\ttfamily,
	  keywordstyle=\bfseries\color{blue},
	  commentstyle=\itshape\color{blue},
	  identifierstyle=\color{blue},
	  stringstyle=\color{orange},
	}
	
	\lstdefinestyle{style2}{
	  belowcaptionskip=1\baselineskip,
	  frame=L,
	  xleftmargin=\parindent,
	  language=C,
	  basicstyle=\footnotesize\ttfamily,
	  commentstyle=\itshape\color{blue},
	}
	\lstset{style=style1}
	
	%FINE STYLE INSERIMENTO CODICE
	%*****************************************
	\usepackage[default]{cantarell} %% Use option "defaultsans" to use cantarell as sans serif only
	\usepackage[T1]{fontenc}        %% for font
	\hypersetup{colorlinks, linkcolor=black, urlcolor=blue}
	\newcommand{\addglos}{\begin{scriptsize}{\textbf{\ped{G}}} \end{scriptsize}} 
	\pagestyle{fancy}
	\fancyhead{}
	\fancyfoot{}
	%\fancyhead[L]{\includegraphics[scale=0.28]{team_not_found.jpeg}}
	\fancyhead[L]{\includegraphics[scale=0.15]{../../team404_small.jpg} \hspace{2mm} QUIZZIPEDIA}
	\fancyhead[R]{\leftmark}
	\fancyfoot[L]{Universit\`a degli studi di Padova - IS 2015/2016 \\ \url{team404swe@gmail.com}}

	
	%Commando usato per la tabella di informazioni sul documento
	\newcommand{\introtab}[9]{
		\begin{table}[ht]
		\begin{center}		
		\begin{tabular}{r l}			
			\toprule		
			\multicolumn{2}{c}{\textbf{ Informazioni sul documento }} \\
			\midrule 
			\textbf{Nome Documento}			& \vline \hspace{3.5 mm} {#1} \\
			\textbf{Versione}				& \vline \hspace{3.5 mm} {#2} \\
			\textbf{Uso} 					& \vline \hspace{3.5 mm} {#3} \\
			\textbf{Data Creazione} 		& \vline \hspace{3.5 mm} {#4} \\
			\textbf{Data Ultima Modifica} 	& \vline \hspace{3.5 mm} {#5} \\
			\textbf{Redazione}				& \vline \hspace{3.5 mm} {#6} \\
											%& \vline \hspace{3.5 mm} {#7} \\	
			\textbf{Verifica} 				& \vline \hspace{3.5 mm} {#7}	\\
			\textbf{Approvazione}			& \vline \hspace{3.5 mm} {#8}\\	
			\textbf{Committente} 			& \vline \hspace{3.5 mm} Zucchetti SPA\\
			\textbf{Lista di distribuzione} & \vline \hspace{3.5 mm} Prof. Vardanega Tullio \\														& \vline \hspace{3.5 mm} TEAM404 \\
	\bottomrule	
	\end{tabular}
	\end{center}
	\end{table}
	}
	% Comando di inizio del registro
	\newcommand{\beginregistro}{
		%\begin{longtable}{{|p{0.10\textwidth}|p{0.20\textwidth}|p{0.15\textwidth}|p{0.50\textwidth}|}}
		\begin{longtable}{{|p{1.5cm}|p{2.5cm}|p{2cm}|p{8cm}|}} 
	 		\hline	
	}
	% commando usato pr inserire una riga al registro delle modifiche
	\newcommand{\rigaregistro}[4]{
		{\footnotesize #1} & {\footnotesize #2} &  {\footnotesize #3} &  {\footnotesize #4} \\
			\hline	
	}
	% Comando di fine registro
	\newcommand{\fineregistro}{ \end{longtable}	}
	
	%************************************************
	% commandi per il GLOSSARIO
	%***********************************************
	% Commando di inizio tabella Glossario
	\newcommand{\beginglos}{
		\begin{longtable}{{p{0.20\textwidth}p{0.65\textwidth}}}	
	}
	% Commando per i termini del glossario
	
	\newcommand{\itemglos}[2]{
		\textbf{#1 :} & {#2} \\ \\ \\
	}
	% Commando fine Glossario
	\newcommand{\fineglos}{ \end{longtable} }
	% Comando per aggiungere una ssezione numerata con lettere al glossario
	\newcommand{\sezione}{
	\subsection{}	
	\rule[0.3pt]{\linewidth}{0.4pt} \\ % Linea orizzontale
	}
	
\newcommand{\sezioneglos}[1] { 
  \newpage
  \cleardoublepage
  \phantomsection
  \addcontentsline{toc}{section}{#1}
  \vspace{11pt}
  \textbf{\huge{#1} } % Lettera grande 
  \\
  \rule[0.3pt]{\linewidth}{0.4pt} \\ % Linea orizzontale
  \fancyhead[R]{#1}
}

\graphicspath{ {./Images/} }
\title{\textbf{{\fontsize{10mm}{6mm}\selectfont QUIZZIPEDIA}}}

\makeindex

\begin{document}
	\maketitle
	
	\begin{center}

	\includegraphics{../../team_not_found.jpg}\\	
	\fontsize{5mm}{3mm}\url{team404swe@gmail.com}\\
	\vspace{40mm}
	\textbf{ Verbale 1 }
	\end{center}
	\thispagestyle{empty}	% per togliere il numero in fondo pagina
	\introtab{Verbale 1}			%1 nome documento
			{1.0} 							%2 versione
			{Esterno} 						%3 Uso
			{8 gennaio 2016} 				%4 Data cre
			{\today} 						%5 Data mod
			{Davide Bortot}					%6 Redazione1
			%{} 							%7 Redazione2
			{Andrea Multineddu} 			%8 Verifica
			{Davide Bortot} 				%9 Approvazione
	
	\newpage
	\section{Sommario}
	Questo documento ha lo scopo di formalizzare l'incontro tra i membri del gruppo \textbf{Team404} avvenuto il 18/12/2015. Scopo dell'incontro era mettere le basi ed accordarsi sui metodi di comunicazione e sulla prima suddivisione dei ruoli tra gli elementi.
	\subsection{Generalità}
	\begin{itemize}
	\item\textbf{Data}: 18 Dicembre 2015.
	\item\textbf{Luogo}: Padova, Via Trieste 63, Dipartimento di Matematica Pura ed Applicata dell'Università di Padova, Aula 1C/150.
	\item\textbf{Ora d'inizio}: 13:00.
	\item\textbf{Durata}: 40 minuti.
	\item\textbf{Partecipanti}: Davide Bortot, Martin Vadice Mbouenda, Marco Crivellaro e Alex Beccaro.
	\end{itemize}

	\newpage
	\section{Argomenti discussi}
		\subsection{Nome del Gruppo}
		Ci sono state alcune proposte come "I giubilati", "I fuorilegge" e "Team404". Quest'ultima è la proposta che ha riscosso maggiore approvazione, quindi "Team404" è da considerarsi nome ufficiale del gruppo.
		\subsection{Mezzi di comunicazione e di gestione}
		Durante l'incontro sono stati proposti i seguenti canali di comunicazione e di collaborazione:
		\begin{itemize}
			\item \textbf{WhatsApp}: adatto per comunicazioni semplici e veloci, sarà il canale ufficiale di comunicazione interna.
			\item \textbf{Dropbox}: verrà creata una cartella condivisa su Dropbox in cui mettere ogni file utile relativo al progetto (documenti, codice, immagini, etc.).
			\item \textbf{E-mail}: creazione di un e-mail (Gmail) di riferimento del gruppo per le comunicazioni esterne ufficiali, principalmente gestita del Responsabile.
			\item \textbf{Skype}: mezzo comunicativo per decisioni importanti o complesse di gruppo in caso non sia possibile trovarsi personalmente. Valutabili anche alternative (TeamSpeak, altre).
			\item \textbf{Gestione di progetto}: sono state proposte tre diverse piattaforme web per la gestione di progetto: Trello, Asana, TeamWork. Serviranno ulteriori discussioni per decidere quale strumento utilizzare.
		\end{itemize}
		\subsection{Documentazione}
		Si è deciso che tutti i documenti dovranno essere creati in LaTeX, per fornire una struttura solida ma allo stesso tempo flessibile alle esigenze dei contenuti. Per assicurare coerenza di stile e formattazione verrà creato un file "template" da includere in ogni documento.
		\subsection{Ruoli}
		Si è deciso che durante il primo periodo del progetto i ruoli verranno così suddivisi:
		\begin{table}[h!]			
		\begin{center}
			\begin{tabular}{l c}
			\textbf{Componente} & \textbf{Ruolo}\\
			\midrule
			Davide Bortot & Responsabile\\
			Martin Vadice Mbouenda & Amministratore\\
			Marco Crivellaro & Analista\\
			Alex Beccaro & Analista\\
			Luca Alessio & Analista/Verificatore\\
			Andrea Multineddu & Analista/Verificatore\\
			\midrule
			\end{tabular}
		\end{center}
		\end{table}
		\newline
		Gli ultimi due membri citati dovranno accordarsi su quale ruolo ricoprire; preferibilmente la figura dell'analista dovrebbe essere disponibile a trovarsi fisicamente con gli altri due analisti.
\end{document}